% from ega4-18-fr.tex
\oldpage[IV-4]{116}

et comme $L_n=L_{n+1}/\mathfrak J^{n+1}L_{n+1}$ et $M_n=M_{n+1}/\mathfrak J^{n+1}M_{n+1}$, les flèches verticales dans le diagramme commutatif
\[
\xymatrix{
0\ar[r] & N_{n+1}\ar[r]^{v_{n+1}}\ar[d] & L_{n+1}\ar[r]^{u_{n+1}}\ar[d] & M_{n+1}\ar[r]\ar[d] & 0\\
0\ar[r] & N_n\ar[r]_{v_n} & L_n\ar[r]_{u_n} & M_n\ar[r] & 0
}
\]
sont toutes trois surjectives. Or on a $M=\varprojlim M_n$ et $L=A^r=\varprojlim L_n$; si l'on pose $N=\varprojlim N_n$, $u=\varprojlim u_n$, $v=\varprojlim v_n$, on a, en vertu de ($0_{\mathrm{III}}$, 13.2.2), la suite exacte
\begin{equation}
  0\longrightarrow N\xrightarrow{v}L\xrightarrow{u}M\longrightarrow 0.
  \tag{18.3.2.2}\label{IV.18.3.2.2}
\end{equation}
En outre, comme pour tout $n$, $v_n$ est inversible à gauche et que $M_n$ est un $A_n$-module projectif, il résulte de (0, 19.1.8) que la suite exacte (18.3.2.2) est scindée, ce qui prouve le lemme.

Cela étant, montrons d'abord que le foncteur de l'énoncé de (18.3.2) est \emph{pleinement fidèle}. Posons, comme dans le lemme, $A_n=A/\mathfrak J^{n+1}$; soient $B$, $C$ deux $A$-algèbres finies et étales, et posons, pour tout $n$, $B_n=B\otimes_A A_n$, $C_n=C\otimes_A A_n$; en vertu de (18.3.1) et (18.3.2.1), $B$ et $C$ sont séparés et complets pour la topologie $\mathfrak J$-pré-adique, et l'on a $B=\varprojlim B_n$, $C=\varprojlim C_n$; en outre tout homomorphisme de $A$-algèbres $u:B\to C$ est continu pour les topologies $\mathfrak J$-pré-adiques, et donne donc un système projectif d'homomorphismes de $A_n$-algèbres $u_n=u\otimes 1:B_n\to C_n$, dont il est la limite projective; la réciproque étant évidente, on a donc une bijection canonique
\[
  \Hom_{A\text{-alg.}}(B,C)\xrightarrow{\sim}\varprojlim\Hom_{A_n\text{-alg.}}(B_n,C_n).
\]
Mais comme $B$ et $C$ sont des $A$-algèbres étales, il résulte aussitôt de (18.1.2) que l'application canonique
\[
  \Hom_{A_{n+1}\text{-alg.}}(B_{n+1},C_{n+1})\longrightarrow\Hom_{A_n\text{-alg.}}(B_n,C_n)
\]
est bijective pour $n\geq 0$, ce qui achève de prouver que l'application canonique
\[
  \Hom_{A\text{-alg.}}(B,C)\longrightarrow\Hom_{A_0\text{-alg.}}(B_0,C_0)
\]
est bijective.

Pour achever de démontrer (18.3.2), il suffit de voir que pour toute $A_0$-algèbre finie et étale $B_0$, il existe une $A$-algèbre finie et étale $B$ et un $A_0$-isomorphisme $B_0\xrightarrow{\sim}B\otimes_A A_0$. Or, il résulte de (18.1.2) qu'il y a un système projectif $(B_n)$ tel que $B_n$ soit une $A_n$-algèbre étale et que les homomorphismes $B_{n+1}\otimes_{A_{n+1}}A_n\to B_n$ soient bijectifs. Il résulte de (18.3.1) et (18.3.2.1) que la $A$-algèbre $B=\varprojlim B_n$ est un $A$-module projectif de type fini et que $B_0$ est isomorphe à $B\otimes_A A_0$. Pour prouver que $B$ est une $A$-algèbre étale, il suffit, en vertu de (18.2.5), de montrer que pour tout idéal maximal $\mathfrak m$ de $A$, $B_{\mathfrak m}/\mathfrak mB_{\mathfrak m}$ est une $(A/\mathfrak m)$-algèbre séparable. Or, puisque $\mathfrak J$ est contenu dans le radical de $A$ ($0_{\mathrm I}$, 7.1.10),

